\documentclass[11pt]{article}
% Shared preamble for per-Python-file documentation (input via \input{...}).
\usepackage[margin=1in]{geometry}
\usepackage[T1]{fontenc}
\usepackage[utf8]{inputenc}
\usepackage{lmodern}
\usepackage{hyperref}
\usepackage{xcolor}
\usepackage{listings}
\usepackage{listingsutf8}
\lstset{
  basicstyle=\ttfamily\small,
  columns=fullflexible,
  breaklines=true,
  upquote=true,
  inputencoding=utf8,
  extendedchars=true,
  frame=single,
  framerule=0.2pt,
  tabsize=2
}


\title{\texttt{run\_tpm.py} (Q\&A)}
\author{}
\date{}

\begin{document}
\maketitle

\paragraph{Q: What is the purpose of \texttt{run\_tpm.py}?}
\textbf{A:} It trains and evaluates the TPM baseline using the tree-based encoding/decoding utilities in \lstinline|utils.py|. It builds percentile-based play-time buckets, trains a TPM classifier over the internal tree decisions, and decodes predictions back to a scalar watch-time estimate.

\paragraph{Q: What are the key hyperparameters unique to TPM training here?}
\textbf{A:} The key parameters are \lstinline|--wr_bucknum| (number of percentile buckets), \lstinline|--variance_weight| (weight on the variance/uncertainty term returned by decoding), \lstinline|--mse_weight| (weight on decoded scalar MSE), and \lstinline|--tree_cla_weight| (weight on the tree classification loss).

\paragraph{Q: How are bucket boundaries determined?}
\textbf{A:} The script uses \lstinline|get_playtime_percentiles_range| to compute percentile boundaries from the training labels, returning \lstinline|bucket_begins| and \lstinline|bucket_ends| tensors used for encoding and decoding.

\paragraph{Q: What loss is optimized during training?}
\textbf{A:} It combines three terms: (1) a tree classification loss over internal nodes (\lstinline|get_tree_classify_loss|), (2) an MSE loss between decoded scalar prediction and the label, and (3) a variance-like regularizer from decoding.

\paragraph{Q: How is a scalar prediction produced at evaluation time?}
\textbf{A:} The test loop calls the model to obtain tree probabilities, decodes them into \lstinline|encoded_y| using \lstinline|get_tree_encoded_value|, and uses that scalar for MAE/XAUC/KL evaluation.

\paragraph{Q: Code example: the combined TPM loss?}
\textbf{A:}
\begin{lstlisting}[language=Python]
loss = (tree_classify_loss * args.tree_cla_weight
        + mse_loss * args.mse_weight
        + variance * args.variance_weight)
\end{lstlisting}

\paragraph{Q: Full source code?}
\textbf{A:}
\lstinputlisting[language=Python]{/workspace/run_tpm.py}

\end{document}
